\chapter{HCRBO/CFI 增量与 unchanged chunk 跳过}\label{ch:hcrbo}

\section{全量 vs 增量：叙事对比}

\begin{analogybox}[图书馆编目比喻]
\textbf{全量 backup}：把一本书从头到尾重新切块、编号、入库。\\
\textbf{增量 backup}：手上有\textbf{上一版编目卡}（CFI 锚点），新书送来时：
\begin{enumerate}[nosep]
  \item 翻到编目卡标注的页码，快速核对「这几页变没变」；
  \item \textbf{没变}：直接复用旧编号（chunk hash），不重新复印、不重新装订；
  \item \textbf{变了}：只对变了的章节重新 FastCDC 切块。
\end{enumerate}
\end{analogybox}

\section{CFI 锚点结构}

每次全量/增量分块可为文件构建 \texttt{CfiIndex}，记录 \texttt{ChunkAnchor}：
\begin{itemize}[nosep]
  \item \texttt{offset, length} — chunk 在文件中的区间；
  \item \texttt{hash[32]} — 上次备份的 Content ID；
  \item \texttt{rolling\_checksum} — 快速预检指纹（v0.1 batch skip）；
  \item \texttt{strength} — Weak（FastCDC 内块）/ Strong（Rabin 边界块）。
\end{itemize}

\begin{figure}[htbp]
\centering
\resizebox{\linewidth}{!}{%
\begin{tikzpicture}[x=0.008cm,y=0.55cm]
  \draw[->] (0,0) -- (12000,0) node[right,font=\small] {文件偏移};
  \foreach \x/\w/\col in {0/800/CoreGreen!40, 800/1200/CoreBlue!30, 2000/900/CoreGreen!40,
                          2900/1100/CoreBlue!30, 4000/1500/CoreGreen!40, 5500/1000/CoreBlue!30,
                          6500/1300/CoreGreen!40, 7800/900/CoreOrange!50, 8700/3300/CoreGreen!40} {
    \fill[\col] (\x,0.4) rectangle (\x+\w,1.3);
  }
  \fill[CoreRed!60] (7800,0.4) rectangle (7900,1.3);
  \node[font=\tiny] at (7850,1.8) {1 字节编辑};
  \foreach \x in {0,800,2000,2900,4000,5500,6500,7800,8700} {
    \draw[thick] (\x,0) -- (\x,1.5);
  }
\end{tikzpicture}%
}
\caption{增量场景：1-byte 编辑后，仅局部区域重切，其余锚点直接 reuse。\\
{\footnotesize 绿色=reuse \quad 蓝色=重切（FastCDC） \quad 橙色=stale 窗口}}
\label{fig:incremental-reuse}
\end{figure}

\section{ChunkIncremental 主循环}

\texttt{EbHcrboChunker::ChunkIncremental}（\filepath{eb\_hcrbo.cc}）按\textbf{排序锚点}扫描文件：

\begin{figure}[htbp]
\centering
\begin{tikzpicture}[node distance=0.6cm, scale=0.95, transform shape]
  \node[block=CoreBlue] (start) {pos=0，sorted\_pos=0};
  \node[decision, below=of start] (ahead) {anchor.offset $>$ pos?};
  \node[block=CoreOrange, below left=1.1cm and 1.0cm of ahead] (gap) {ChunkRegion(pos$\rightarrow$anchor)：\\新区域 FastCDC};
  \node[decision, below right=1.1cm and 1.0cm of ahead] (verify) {rolling + Bloom + SHA256};
  \node[block=CoreGreen, below=of verify] (reuse) {输出 desc，\\reused\_from\_cfi=true};
  \node[block=CoreRed, below left=0.9cm and 0.5cm of verify] (stale) {rolling 不匹配：\\扩展 stale 窗口};
  \node[block=CorePurple, below right=0.9cm and 0.5cm of verify] (recut) {hash 不匹配：\\ChunkRegion 重切};
  \draw[arr] (start) -- (ahead);
  \draw[arr] (ahead) -- node[lab,left]{是} (gap);
  \draw[arr] (ahead) -- node[lab,right]{对齐} (verify);
  \draw[arr] (verify) -- node[lab,right]{匹配} (reuse);
  \draw[arr] (verify) -- node[lab,left]{rolling$\neq$} (stale);
  \draw[arr] (verify) -- node[lab,right]{hash$\neq$} (recut);
\end{tikzpicture}
\caption{增量分块决策流程（简化）。}
\label{fig:hcrbo-loop}
\end{figure}

\section{三级验证：如何跳过 unchanged chunk}

\subsection{第 1 级：Rolling Checksum 预检}

若锚点带 \texttt{rolling\_checksum} $\neq 0$：
\[
\text{rc}_{\text{now}} = \texttt{CfiRollingVerifier::VerifyAnchor}(\text{data},\,offset,\,length)
\]
\begin{itemize}[nosep]
  \item $\text{rc}_{\text{now}} = \text{rc}_{\text{hist}}$ $\Rightarrow$ 可能未变，进入第 2 级；
  \item 不等 $\Rightarrow$ \textbf{rolling mismatch}，batch skip 向前扫描连续 stale 锚点，扩展重切窗口（避免逐锚点 SHA256）。
\end{itemize}

\subsection{第 2 级：CFI Bloom 负向预滤}

rolling 匹配后，查 Bloom filter：
\begin{itemize}[nosep]
  \item \texttt{MightContain(hash)} = false $\Rightarrow$ \textbf{必定变更}，跳过 SHA256（\texttt{cfi\_bloom\_skip\_hits}++）；
  \item true $\Rightarrow$ 进入第 3 级完整 hash。
\end{itemize}

\begin{analogybox}[Bloom = 快速黑名单]
「这本书页码看起来一样，但 Bloom 说上次 hash 绝不可能在这」$\Rightarrow$ 直接判变更，不必算 SHA256。
\end{analogybox}

\subsection{第 3 级：SHA256 内容比对}

\[
\texttt{ContentHash}(\text{data}+offset,\,length) \stackrel{?}{=} \texttt{anchor.hash}
\]
匹配 $\Rightarrow$ 输出 \texttt{ChunkDescriptor}，设 \texttt{reused\_from\_cfi=true}，\textbf{不调用 FastCDC}。

\section{Store 阶段：reuse 如何跳过 IO}

分块阶段标记 \texttt{reused\_from\_cfi} 后，Pipeline Store 路径：
\begin{enumerate}[nosep]
  \item \texttt{Exists(hash)} 检查仓库是否已有该 Content ID；
  \item 已存在 $\Rightarrow$ \texttt{skip\_store=true}，不写 EbPack；
  \item manifest 仍记录该 hash（引用已有 chunk）。
\end{enumerate}

\begin{table}[htbp]
\centering
\caption{EbHcrboStats 关键计数器}
\label{tab:hcrbo-stats}
\small
\begin{tabular}{@{}lp{9cm}@{}}
\toprule
字段 & 含义 \\
\midrule
\texttt{chunks\_reused\_from\_cfi} & rolling+hash 验证通过，直接 reuse \\
\texttt{cfi\_rolling\_skip\_hits} & rolling 不匹配时 batch 跳过的锚点数 \\
\texttt{cfi\_bloom\_skip\_hits} & Bloom 否定后省去 SHA256 的次数 \\
\texttt{chunks\_cut\_fastcdc} & 变更区域内新 FastCDC 切出的块数 \\
\bottomrule
\end{tabular}
\end{table}

\section{HCRBO 与 FastCDC 的关系}

\begin{corebox}[分层不替代]
HCRBO \textbf{不替换} FastCDC 算法；它在 FastCDC 之上提供\textbf{增量路由}：
\begin{itemize}[nosep]
  \item 未变区域 $\rightarrow$ reuse 旧 anchor（跳过 CDC+digest+encode+store）；
  \item 变更区域 $\rightarrow$ 内嵌 \texttt{FastCdcSlice} 重新 \texttt{ChunkRegion}；
  \item Rabin 强锚点（可选）用于区域边界稳定切分。
\end{itemize}
\textbf{Content ID 语义不变}：reuse 的 chunk 仍引用同一 SHA256。
\end{corebox}

\section{Bench 约束}

L2 增量门禁：在 5\,MB 文件 offset 5\,MB 处做 \textbf{1-byte 编辑} 后，reuse $\geq 90\%$（immutable floor）。\\
证明 CFI 路径在典型小编辑下\textbf{绝大多数 chunk 被跳过}。
